\documentclass[11pt,a4paper]{article}

\usepackage[a4paper,margin=2.3cm]{geometry}
\usepackage{fontspec}
\usepackage{titlesec}
\usepackage{longtable}
\usepackage{booktabs}
\usepackage{array}
\usepackage{xcolor}
\usepackage[most]{tcolorbox}
\usepackage{hyperref}
\usepackage{enumitem}
\usepackage{fancyhdr}
\usepackage{amssymb}

\setmainfont{Sylfaen}

\definecolor{brandgreen}{HTML}{1F6F5C}
\definecolor{quoteblue}{HTML}{1D4E89}
\definecolor{termorange}{HTML}{B5651D}
\definecolor{importred}{HTML}{B03A2E}

\hypersetup{
  colorlinks=true,
  linkcolor=brandgreen,
  urlcolor=brandgreen,
  pdftitle={Clarke et al. 2024 (OpenDengue) -- ქართული თარგმანი და შენიშვნები},
  pdfauthor={სერგი კონიაშვილი}
}

\titleformat{\section}{\large\bfseries\color{brandgreen}}{\thesection}{0.5em}{}
\titleformat{\subsection}{\normalsize\bfseries}{\thesubsection}{0.5em}{}

\setlength{\headheight}{14pt}
\pagestyle{fancy}
\fancyhf{}
\fancyhead[L]{Clarke et al. (2024), Scientific Data}
\fancyhead[R]{ქართული თარგმანი სასწავლად}
\fancyfoot[C]{\thepage}
\setlength{\emergencystretch}{3em}
\setlength{\parskip}{0.4em}

\newtcolorbox{etl}{enhanced,breakable,colback=brandgreen!6,colframe=brandgreen,
  fonttitle=\bfseries,title={ETL: მნიშვნელოვანია ჩემი სისტემისთვის}}
\newtcolorbox{citat}{enhanced,breakable,colback=quoteblue!5,colframe=quoteblue,
  fonttitle=\bfseries,title={ციტატა თეზისისთვის (ორიგინალი ინგლისურად)}}
\newtcolorbox{termin}{enhanced,breakable,colback=termorange!7,colframe=termorange,
  fonttitle=\bfseries,title={ტერმინები}}
\newtcolorbox{skipbox}{enhanced,breakable,colback=gray!7,colframe=gray!55,
  fonttitle=\bfseries,title={ნაკლებად მნიშვნელოვანია შენთვის: საკმარისია ზოგადი გაგება}}
\newtcolorbox{alarm}{enhanced,breakable,colback=importred!5,colframe=importred,
  fonttitle=\bfseries,title={ყურადღება: ეს შენს სისტემას პირდაპირ ეხება}}
\newcommand{\imp}[1]{\textbf{\color{importred}#1}}
\newcommand{\pg}[1]{{\small\color{gray}[ორიგინალი PDF, გვ.~#1]}}

\begin{document}

\begin{center}
  {\Large\bfseries A global dataset of publicly available dengue case count data}\\[0.6em]
  {\large დენგეს შემთხვევების საჯაროდ ხელმისაწვდომი მონაცემების გლობალური ნაკრები}\\[0.8em]
  Clarke J., Lim A., Gupte P., Pigott D.M., van Panhuis W.G., Brady O.J.\\
  \textit{Scientific Data} 11, 296 (2024), ტიპი: \textit{Data Descriptor}\\
  DOI: \href{https://doi.org/10.1038/s41597-024-03120-7}{10.1038/s41597-024-03120-7}\\
  მონაცემთა ნაკრები (Figshare): \href{https://doi.org/10.6084/m9.figshare.24259573}{10.6084/m9.figshare.24259573}
\end{center}

\noindent\rule{\linewidth}{0.4pt}

\noindent\textbf{ეს დოკუმენტი სასწავლო თარგმანია და ოფიციალური არ არის.} სტატია ღიაა, ამიტომ წყაროს მითითებით თარგმნა დასაშვებია.
თეზისში \imp{ყოველთვის ორიგინალი ინგლისური ტექსტი} ციტირე. სრულადაა თარგმნილი: Abstract, Background \& Summary,
Usage Notes, Citation \& licence. Methods და Technical Validation ნაწილობრივ შემოკლებულია, ეს ადგილები \textit{[შემოკლებით]} არის მონიშნული.

\subsection*{რატომ არის ეს სტატია შენთვის მნიშვნელოვანი}
ეს \textbf{მონაცემთა სტატიაა} (Data Descriptor) და არა მოდელირების: ავტორები აღწერენ, როგორ შექმნეს OpenDengue-ს მონაცემთა ბაზა.
შენი სისტემა დენგეს მონაცემებს \imp{სწორედ OpenDengue-დან} იღებს. სტატია შენთვის სამ რამეს იძლევა:
\begin{enumerate}[leftmargin=1.5em]
  \item \textbf{მოტივაცია}: ავტორები თვითონ წერენ, რომ მონაცემების პოვნა, ამოღება და სტანდარტიზაცია დიდ დროს მოითხოვს.
  \item \textbf{შენი წყაროს გაგება}: რა არის OpenDengue-ში, რა შეზღუდვები აქვს და როგორ უნდა იქნას გამოყენებული.
  \item \textbf{მოთხოვნები შენს სისტემას}: სწორი ამონაწერის (extract) არჩევა, ვერსიის მითითება, ლიცენზია (CC BY-SA!), კვირეული მონაცემები.
\end{enumerate}

\subsection*{ფერების ლეგენდა}
\textcolor{brandgreen}{\textbf{მწვანე}}: მნიშვნელოვანია შენი სისტემისთვის;
\textcolor{quoteblue}{\textbf{ლურჯი}}: ციტატა თეზისისთვის;
\textcolor{termorange}{\textbf{ნარინჯისფერი}}: ტერმინები;
\textcolor{importred}{\textbf{წითელი ყუთი}}: შენი სისტემის პრობლემა ან შესამოწმებელი საკითხი;
\textcolor{gray}{\textbf{ნაცრისფერი}}: ნაკლებად მნიშვნელოვანი.

%==========================================================
\section{Abstract (რეზიუმე)} \pg{1}

OpenDengue არის დენგეს შემთხვევების მონაცემების გლობალური ბაზა, \imp{საჯარო წყაროებიდან შეგროვებული, სტანდარტიზებული
და ისე ფორმატირებული, რომ ხელახალი ანალიზი გაადვილდეს}. ბაზის 1.2 ვერსია 102 ქვეყნიდან 56 მილიონზე მეტ დენგეს შემთხვევას
მოიცავს, 1924--2023 წლებისთვის, რაც მას ამჟამად ყველაზე დიდ და ყოვლისმომცველ დენგეს მონაცემთა ბაზად აქცევს.
\imp{ჩანაწერების 95\%-ზე მეტი კვირეული ან თვიური რეზოლუციისაა}, ხოლო \imp{სუბნაციონალური მონაცემები 40 ქვეყნისთვის} არის
ხელმისაწვდომი. OpenDengue-ს შესაქმნელად სისტემატურად მოვიძიეთ მონაცემთა ბაზები, ჯანდაცვის სამინისტროების ვებგვერდები,
რეცენზირებული ლიტერატურა და ProMED-mail-ის ანგარიშები, და ამოვიღეთ მნიშვნელზე დაფუძნებული (denominator-based) შემთხვევების
მონაცემები. თანმიმდევრულობის უზრუნველსაყოფად და შეუსაბამობების მოსაგვარებლად სტანდარტიზაციისა და შეცდომების შემოწმების
პროტოკოლებს ვიყენებთ. \imp{ამოღების პროცესი დეტალურად დავადოკუმენტირეთ, რათა ჩანაწერები წყაროსთან დაკავშირებადი და
რეპროდუცირებადი იყოს.} OpenDengue-ს ბაზა კვლავ ვითარდება, დაგეგმილია შემდგომი დეტალიზაცია, და მომხმარებლების წვლილი
მისასალმებელია. ეს ახალი მონაცემთა ნაკრები შეიძლება გამოყენებულ იქნას დენგეს გადაცემის გრძელვადიანი მამოძრავებელი
ფაქტორების უკეთ გასაგებად, დაავადების ტვირთის შეფასების გასაუმჯობესებლად, ინტერვენციების დაგეგმვისა და შეფასებისთვის და
მომავლის პროგნოზების გასაუმჯობესებლად.

\begin{termin}
\textbf{Data Descriptor}: სტატიის ტიპი, რომელიც მონაცემთა ნაკრებს აღწერს (როგორ შეიქმნა, რას შეიცავს), ახალ მეცნიერულ აღმოჩენას კი არა.\\
\textbf{Denominator-based}: მონაცემი, რომელსაც მკაფიო მნიშვნელი აქვს: განსაზღვრული მოსახლეობა, ადგილი და დრო (მაგ. „10 შემთხვევა ქალაქ X-ში, 2015 წლის იანვარში“).\\
\textbf{Subnational / Admin1 / Admin2}: ქვეყნის ადმინისტრაციული დაყოფა. Admin0 = ქვეყანა, Admin1 = პროვინცია ან რეგიონი, Admin2 = რაიონი.\\
\textbf{Passive surveillance}: შემთხვევები აღირიცხება მხოლოდ მაშინ, როცა ავადმყოფი თვითონ მიდის ექიმთან.
\end{termin}

%==========================================================
\section{Background \& Summary (ფონი და შეჯამება)} \pg{1--2}

დენგე გლობალური საზოგადოებრივი ჯანდაცვისთვის მნიშვნელოვანი, მზარდი ინფექციური დაავადებაა: შეფასებით, წელიწადში 100 მილიონი
სიმპტომური ინფექცია 125-ზე მეტ ქვეყანაში. დენგეს ვირუსს (DENV) \textit{Aedes} კოღოები ავრცელებენ და ის წელიწადში 10\,000
სიკვდილს იწვევს. \imp{დენგეს გადაცემისთვის ხელსაყრელი გარემო ფართოვდება კლიმატის ცვლილების}, ურბანიზაციისა და საერთაშორისო
მოგზაურობის გამო. პროგნოზით, 2080 წელს 2015 წელთან შედარებით 2.25 მილიარდით მეტი ადამიანი იქნება დენგეს რისკის ქვეშ,
სულ 6.1 მილიარდი, ანუ მსოფლიოს მოსახლეობის 60\%-ზე მეტი.

\imp{დენგეს ტვირთის გავრცელების თვალყურის დევნება რთულია, რადგან თანმიმდევრული და შესადარებელი მონაცემების შეგროვება და
აგრეგირება რთულია.} ყველაზე ხშირად ხელმისაწვდომი საზომი პასიური ზედამხედველობის შემთხვევებია: ადამიანი სიმპტომებით
მიმართავს სამედიცინო დაწესებულებას, სადაც კლინიკური ალგორითმებით ან ლაბორატორიული დიაგნოსტიკით დგინდება საეჭვო, სავარაუდო ან
დადასტურებული შემთხვევა. შემდეგ ეს მონაცემები რამდენიმე დამუშავების ეტაპს გადის, ჩვეულებრივ ადგილობრივ ან რეგიონულ ჯანდაცვის
დეპარტამენტებსა და ეროვნულ სამინისტროებში. სამინისტროები აგრეგირებულ სტატისტიკას სხვადასხვა სისრულით აქვეყნებენ
ეპიდემიოლოგიურ ბიულეტენებში, ანგარიშებსა და დაფებზე (dashboards).

\begin{citat}
\textit{„Tracking the expansion of the burden of dengue is challenging due to the difficulties in collecting and aggregating consistent and comparable dengue incidence and prevalence data.“} (Background \& Summary, p.~1)
\end{citat}

\imp{მიუხედავად იმისა, რომ ბევრი ქვეყანა დენგეს სტატისტიკას რეგულარულად აქვეყნებს, ამ მონაცემების პოვნა ხშირად რთულია და არ
არსებობს ერთიანი ბაზა, რომელიც სხვადასხვა ქვეყნის მონაცემებს აერთიანებს} გლობალური ტენდენციების შესაფასებლად. ყველა ენდემური
რეგიონის მონაცემების შეგროვება საშუალებას მისცემდა ხელახალ ანალიზს, რათა \imp{უკეთ გავიგოთ გადაცემის მამოძრავებელი ფაქტორები},
ვაკონტროლოთ დაავადების შემცირების მიზნების მიღწევა, შევაფასოთ ინტერვენციები და \imp{მოვახდინოთ მომავალი ტვირთის მოდელირება
სხვადასხვა კლიმატურ სცენარში}. რაც უფრო მაღალია მონაცემების სივრცითი და დროითი რეზოლუცია, მით უფრო ინფორმატიული და ადგილობრივად
სპეციფიკურია ანალიზი.

\begin{citat}
\textit{„While many countries regularly publish dengue case statistics, they can often be difficult to find and no single database aggregates data from multiple countries to assess trends at the global level. Gathering data across all dengue endemic regions would enable re-analysis to better understand the drivers of transmission, [\ldots] and model the possible future burden and spatial limits under different climate scenarios.“} (Background \& Summary, p.~1)
\end{citat}

\textit{[შემოკლებით]} რეგიონული და გლობალური ბაზების შექმნის მცდელობები არსებობდა, მაგრამ თითოეულს შეზღუდვები ჰქონდა (Table 1):
\begin{itemize}[leftmargin=1.5em]
  \item \textbf{WHO}: ქვეყნის დონის მონაცემებს წელიწადში ერთხელ ან ორჯერ იღებს, რაც დეტალური ანალიზისთვის საკმარისად ხშირი არ არის.
  \item \textbf{Project Tycho}: 80 ქვეყანა, 1960--2012, მანქანურად წაკითხვადი, მაგრამ 2012 წლის შემდეგ აღარ განახლებულა.
  \item \textbf{PAHO PLISA}: კვირეული მონაცემები, კარგი მეტამონაცემები, მაგრამ მხოლოდ ამერიკის კონტინენტზე. ორიენტირებულია კუმულაციურ რაოდენობებზე.
  \item \textbf{DengueNet / Dengue Explorer}: თანმიმდევრულობის პრობლემები, მხოლოდ ქვეყნის დონე.
  \item \textbf{ECDC}: მხოლოდ ევროპა (არაენდემური), წლიური.
  \item \textbf{Global Burden of Disease}: ეროვნული შეფასებები, მაგრამ ამოღებული ცხრილები საჯარო არ არის.
  \item \textbf{GIDEON}: ფასიანი გამოწერაა.
  \item \textbf{ProMED-mail}: ყოველდღიური ანგარიშები, მაგრამ მანქანურად წაკითხვადი ფორმატის გარეშე, და \imp{„საჭიროებს მნიშვნელოვან ხელით დამუშავებას“}.
\end{itemize}

\begin{citat}
\textit{„To date, no repository has been able to combine global coverage, public availability, machine readable accessible formats at high spatial-temporal resolution and sustained updates over long time periods.“} (Background \& Summary, p.~2)
\end{citat}

დღემდე ვერცერთმა საცავმა ვერ შეძლო ერთდროულად გლობალური მოცვა, საჯარო ხელმისაწვდომობა, მანქანურად წაკითხვადი ფორმატები მაღალი
სივრცით-დროითი რეზოლუციით და ხანგრძლივი პერიოდის განმავლობაში მუდმივი განახლება. თუმცა ორმა ბოლოდროინდელმა მოვლენამ ეს ამოცანა
დენგესთვის უფრო მისაღწევი გახადა. პირველი: ჯანდაცვის ზედამხედველობის ციფრულ სისტემებში (მაგ. DHIS2) მსოფლიო ინვესტიციებმა
მონაცემების მოცვა, სიჩქარე, საიმედოობა და ხელმისაწვდომობა გაზარდა. მეორე: COVID-19 პანდემიამ აჩვენა, რამდენად საჭიროა
ინფექციური დაავადებების მონაცემების საჯაროდ ხელმისაწვდომობა.

\imp{დენგეს შემთხვევების მონაცემები მრავალ ფორმატში და მრავალი წყაროდან მოდის, რაც დამუშავების სხვადასხვა მეთოდოლოგიას
მოითხოვს.} დეტალური მეტამონაცემები მნიშვნელოვანია, რათა რიცხვები ორიგინალ წყარომდე იყოს მიკვლევადი და წყაროების შედარებადობა
შეფასდეს. \imp{მონაცემების მოძიება, ამოღება, დამუშავება და სტანდარტიზაცია დროს მოითხოვს, მაგრამ აუცილებელია ხელახალი გამოყენებისა
და ანალიზისთვის.} აქ აღვწერთ, როგორ მოვიძიეთ, ამოვიღეთ და დავაფორმატეთ დენგეს დროითი სერიები ენდემური სამყაროდან მაქსიმალური
სივრცითი და დროითი რეზოლუციით, და როგორ შეიფუთა ეს მონაცემები საჯარო ბაზასა და ვებგვერდზე.

\begin{citat}
\textit{„Dengue case data exists in multiple formats from a wide variety of sources that require various processing methodologies. Detailed source metadata is important to ensure case counts can be traced back to their original reporting source, and to enable assessment of comparability between sources. Locating, extracting, processing and standardising data all takes time but is essential to enable reuse and re-analysis.“} (Background \& Summary, p.~2)
\end{citat}

\begin{etl}
ეს სამი აბზაცი \imp{შენი მოტივაციის ცენტრალური ციტატებია}. ისინი \textbf{ორ რამეს} ამტკიცებს:
\begin{enumerate}[leftmargin=1.5em]
  \item \textbf{მონაცემების მომზადება შრომატევადია} („Locating, extracting, processing and standardising data all takes time“). ეს ავტორების საკუთარი სიტყვებია და Suttirat-ისგან განსხვავებით აქ ეს \textbf{პირდაპირ წერია}!
  \item \textbf{წყაროსთან მიკვლევადობა მნიშვნელოვანია} („traced back to their original reporting source“). ეს შენს provenance / audit trail ფუნქციას (FR-18) ასაბუთებს.
\end{enumerate}
\textbf{ფრთხილად:} ეს ციტატები \textbf{დენგეს შემთხვევების} მონაცემებს ეხება, ანუ ჯანმრთელობის მხარეს. OpenDengue ამ პრობლემას
ჯანმრთელობის მხრიდან წყვეტს. \textbf{კლიმატის მონაცემებთან გაერთიანება} კი შემდეგი ნაბიჯია, რომელიც OpenDengue-ს არ აქვს, მაგრამ
ავტორები მას გამოყენების მიზნად ასახელებენ („drivers of transmission“, „climate scenarios“). ამ ნაბიჯს შენი სისტემა აკეთებს.
ეს ლოგიკური კავშირი \textbf{შენი არგუმენტია} და თეზისში ასეც უნდა ჩამოაყალიბო.
\end{etl}

%==========================================================
\section{Methods (მეთოდები)} \pg{2--7}

\begin{skipbox}
Methods აღწერს, როგორ შეაგროვეს ავტორებმა მონაცემები. დეტალები შენთვის ნაკლებად მნიშვნელოვანია, მაგრამ \textbf{ზოგადი სურათი}
ძალიან კარგი მაგალითია იმისა, \textbf{რამდენად შრომატევადია} ხელით ETL. ქვემოთ მოცემულია მოკლე თარგმანი და შენიშვნები.
\end{skipbox}

\subsection*{საძიებო სტრატეგია \textit{[შემოკლებით]}} \pg{2--3}
ოთხი ტიპის წყარო მოიძიეს: ჯანდაცვის სამინისტროების ვებგვერდები, არსებული ბაზები, რეცენზირებული სტატიები და ProMED-mail.
შექმნეს წყაროების პრიორიტეტების იერარქია. ბიულეტენები და წლიური ანგარიშები \imp{საიტის რუკებითა და ძიების ფუნქციებით
მოიძებნა}. \imp{ინგლისური ენის გარეშე ვებგვერდები Google Translate-ით} და ამ ქვეყნების კოლეგების დახმარებით დაათვალიერეს.
სტატიები PubMed-ით მოიძებნა. ძიება უფრო მიზანმიმართული და ეფექტური გახდა, \imp{რაც უფრო კარგად გაიცნეს თითოეული ქვეყნის
ანგარიშგების სისტემა}. მონაცემთა ხარვეზები რუკებით (heatmaps) ფასდებოდა და მიზანმიმართულად ივსებოდა.
შერჩევის კრიტერიუმები: წყარო \imp{საჯარო} და \imp{მნიშვნელზე დაფუძნებული} უნდა იყოს. შემოტანილი (იმპორტირებული)
შემთხვევები გამოირიცხა, სადაც ცალკე იყო მოცემული.

\subsection*{მონაცემების ამოღება \textit{[შემოკლებით]}} \pg{3--4}
ფაილები სხვადასხვა ფორმატში იყო: \imp{.csv, .XLS და .pdf}. სამინისტროების მონაცემები ჩამოიტვირთა როგორც \imp{გრაფიკები, ცხრილები და
ტექსტი}. PAHO PLISA-დან 53 ცალკეული კვირეული ჩამოტვირთვა გაკეთდა 56 ქვეყნისთვის. ყველა ტექსტი UTF-8-ში გადაიყვანეს.
თითოეულ წყაროს ფაილს მიენიჭა უნიკალური იდენტიფიკატორი (UUID, მაგ. \texttt{MOH-MEX-2012-Y01-01}).
\begin{itemize}[leftmargin=1.5em]
  \item ცხრილი PDF-ში $\rightarrow$ ამოღებული R პაკეტებით (\texttt{pdftools}, \texttt{tabulizer}) ან Excel-ით;
  \item გრაფიკი $\rightarrow$ რიცხვები \imp{WebPlotDigitizer-ით} ამოიღეს (გრაფიკიდან წერტილების „წაკითხვა“);
  \item ტექსტი $\rightarrow$ Google Translate-ით თარგმნეს და \imp{ხელით გადაიტანეს ცხრილში};
  \item რუკა $\rightarrow$ სადაც შესაძლებელი იყო, \imp{ხელით} ამოიღეს;
  \item ყველაფერი შემდეგ R-ით დამუშავდა, \imp{სვეტების სახელები გასწორდა} და ერთ \texttt{.csv} ფაილად გაერთიანდა.
\end{itemize}

\begin{etl}
ეს სია ხელით ETL-ის შრომატევადობის \imp{შესანიშნავი ილუსტრაციაა}: PDF-ების სკრეიპინგი, გრაფიკებიდან რიცხვების ამოღება, თარგმნა, ხელით
გადაწერა, სვეტების სტანდარტიზაცია. ოღონდ ეს OpenDengue-ს გუნდმა \textbf{უკვე გააკეთა}. შენი სისტემა ამ სამუშაოს
შედეგს იყენებს და მას კლიმატს უმატებს. თეზისში შეგიძლია დაწერო: „ზოგი წყარო (როგორიცაა OpenDengue) ჯანმრთელობის მონაცემებს
უკვე ამზადებს, მაგრამ მისი კლიმატის მონაცემებთან გაერთიანება კვლავ მკვლევრის საქმეა“.
\end{etl}

\subsection*{წყაროს მონაცემები და მეტამონაცემები \textit{[შემოკლებით]}} \pg{4}
ყოველი წყაროსთვის შეინახა: \imp{წვდომის თარიღი, მთავარი გვერდის URL, ვებგვერდზე ნავიგაციის ნაბიჯები, საძიებო სიტყვები} და
შენიშვნები. ჩანაწერის UUID-ით შეიძლება წყაროს პოვნა \texttt{sourcedata\_V1.2.csv}-ში. შემთხვევის განსაზღვრებები სამ დონეზე
სტანდარტიზდა: \textit{suspected} (საეჭვო), \textit{probable} (სავარაუდო), \textit{confirmed} (დადასტურებული). მეტამონაცემები
\texttt{.json} ფორმატშია, NIAID Data Ecosystem-ის სქემით (schema.org-ზე დაფუძნებით).

\begin{etl}
\textbf{მეტამონაცემების ეს დონე} (თარიღი, URL, ნაბიჯები) კარგი მოდელია შენი provenance-ისთვის. შეადარე: რას იწერს შენი სისტემა
ყოველ წყაროზე? მაგალითად, \imp{ჩამოტვირთვის თარიღს} და \imp{მონაცემთა ნაკრების ვერსიას} თუ იწერს? თუ არა, ეს რეკომენდაცია A5-ის
(რეპროდუცირებადობის პაკეტი) დამატებითი დასაბუთებაა.
\end{etl}

\subsection*{მონაცემთა დამუშავება \textit{[შემოკლებით]}} \pg{4--5}
\textbf{გეო-შესაბამისობა (Geomatching).} ადგილის სახელები ორ საერთაშორისო shapefile ფორმატს დაუკავშირდა: FAO GAUL (Admin0--2)
და Natural Earth (Admin0--1). ქვეყნები ISO alpha-3 კოდებს დაუკავშირდა R პაკეტით \texttt{countrycode}, ქვე-ერთეულები კი
\imp{იერარქიული „fuzzy“ შესაბამისობით} (\texttt{hmatch}). რაც ვერ დაემთხვა, \imp{ხელით გასწორდა}, საჭიროების შემთხვევაში
Google Maps-იდან კოორდინატების აღებით.

\textbf{დროითი შესაბამისობა (Time matching).} თარიღები კალენდარულ ფორმატში გადაიყვანეს. ზოგი წყარო \imp{„ეპიდემიოლოგიურ კვირებს“}
იყენებდა. ეს კალენდარულ დაწყებისა და დასრულების თარიღებად გადაიყვანეს R პაკეტით \texttt{EpiWeek}: \imp{წლის პირველი ეპი-კვირა
იანვრის სულ მცირე 4 დღეს შეიცავს, კვირა იწყება კვირა დღეს და მთავრდება შაბათს} (აშშ-ის CDC-ის ვერსია).

\begin{alarm}
\textbf{კვირეული მონაცემები $\rightarrow$ თვიური:} OpenDengue-ს ჩანაწერებს აქვს \texttt{calendar\_start\_date} და \texttt{calendar\_end\_date}.
ეპი-კვირა შეიძლება \imp{ორ თვეს შორის გადადიოდეს} (მაგ. 29 იანვარი -- 4 თებერვალი). შეამოწმე: როგორ შეიქმნა შენი
\texttt{dengue\_monthly\_by\_iso3.csv}? რომელ თვეს მიეკუთვნა ასეთი კვირა? ეს ETL-ის მნიშვნელოვანი გადაწყვეტილებაა და თეზისში
უნდა აღწერო. თუ სისტემაში კვირეულ რეზოლუციას დაამატებ (A4), ეპი-კვირის ეს განმარტება გამოიყენე.
\end{alarm}

\textbf{ურთიერთსაწინააღმდეგო ჩანაწერები.} თუ ერთი ადგილისა და დროისთვის რამდენიმე წყარო სხვადასხვა რიცხვს აჩვენებდა: ზუსტი
დუბლიკატები წაიშალა; ერთი კატეგორიის წყაროებიდან უდიდესი მნიშვნელობა აირჩა; სხვადასხვა კატეგორიიდან კი პრიორიტეტით:
სამინისტრო > რეგიონული ორგანო > Project Tycho > სტატიები > სხვა.

\imp{ამიტომ OpenDengue სამ ამონაწერს (extract) გვთავაზობს:} (i) საუკეთესო ეროვნული (Admin0) შეფასება, (ii) მაქსიმალური დროითი
დეტალიზაცია, (iii) მაქსიმალური სივრცითი დეტალიზაცია. ავტორებმა ჩანაწერები \imp{სივრცით შესაბამისობაში არ მოიყვანეს}: Admin2-ის
ჯამი შეიძლება Admin0-ის რიცხვს არ ემთხვეოდეს.

\begin{alarm}
\textbf{რომელ ამონაწერს იყენებს შენი სისტემა?} Usage Notes-ში (ქვემოთ) ავტორები პირდაპირ წერენ, რომ \imp{დროითი სერიების ანალიზისა და
პროგნოზირებისთვის \textbf{temporal extract}} უნდა იქნას გამოყენებული. შენი სისტემა სწორედ დროით სერიებს ამზადებს. შეამოწმე,
რომელი ამონაწერიდან და რომელი ვერსიიდან შეიქმნა შენი ფაილი, და ეს სისტემის დოკუმენტაციაში ჩაწერე. თუ სუბნაციონალურ ფუნქციას (A8)
დაამატებ: Admin1-ის ჯამი შეიძლება ეროვნულ რიცხვს არ ემთხვეოდეს, ასე რომ სისტემამ ეს \textbf{არ უნდა „გაასწოროს“} და
მომხმარებელს უნდა აუხსნას.
\end{alarm}

\textbf{დენგეს კლასიფიკაცია.} სვეტი \texttt{dengue\_total} ყველა სიმძიმის (დენგე, მძიმე დენგე, DHF, DSS, სიკვდილი) და დადასტურების
ყველა მეთოდის შემთხვევებს მოიცავს. სტანდარტიზებული განსაზღვრება სვეტ \texttt{case\_definition\_standardised}-შია.

\begin{etl}
\textbf{Suttirat-თან განსხვავება:} Suttirat-ი მხოლოდ \textbf{DHF}-ს (მძიმე ფორმას) იყენებს, OpenDengue-ს \texttt{dengue\_total} კი
\textbf{ყველა} შემთხვევას მოიცავს. ამიტომ Case Study 1-ში შენი სისტემის რიცხვები Suttirat-ის რიცხვებს \imp{ვერ დაემთხვევა} და ეს
ნორმალურია. თეზისში ეს განსხვავება ახსენი.
\end{etl}

\subsection*{PAHO PLISA: კუმულაციურიდან ინციდენტურზე \textit{[შემოკლებით]}} \pg{5--7}
\begin{skipbox}
ეს ნაწილი მხოლოდ ამერიკის კონტინენტის მონაცემებს ეხება. PAHO კუმულაციურ რიცხვებს აქვეყნებს („29 იანვრამდე სულ X შემთხვევა“),
ავტორებს კი კვირეული რიცხვები სჭირდებოდათ. ოთხი პრობლემა გადაჭრეს: მოულოდნელი ნახტომები, მნიშვნელობების შემცირება გადასინჯვისას
(უარყოფითი კვირეული რიცხვები $\rightarrow$ NA), ხარვეზები და მონაცემების არარსებობა. გამოიყენეს \textbf{„zero filling“} (წლიური ჯამის დახმარებით
ნულებით შევსება) და 6 კვირამდე ხარვეზების \textbf{cubic spline ინტერპოლაცია} (R პაკეტი \texttt{zoo}). ასეთი ჩანაწერები UUID-ში
\texttt{(Zero filling)} ან \texttt{(Imputed)} სუფიქსით არის მონიშნული, რათა მომხმარებელმა სურვილისამებრ ამოიღოს.
\end{skipbox}

\begin{etl}
ეს ორ რამეს გასწავლის: (1) \imp{იმპუტირებული მონაცემები უნდა იყოს მონიშნული}, რომ მომხმარებელმა იცოდეს, რომელი რიცხვია
ნამდვილი და რომელი შევსებული. შენი სისტემა mass-preserving downscaling-ის დროს (FR-20) ხომ მონიშნავს გამოთვლილ მნიშვნელობებს?
(2) ავტორები \imp{ჯამის შენარჩუნებას} ამოწმებენ (კვირეული ჯამი = წლიური ჯამი). ეს ზუსტად შენი FR-20-ის იდეაა, ამიტომ შეგიძლია
ციტირო, როგორც მსგავსი მიდგომის მაგალითი.
\end{etl}

\subsection*{ვებგვერდი OpenDengue.org} \pg{7}
ვებგვერდი R Markdown-ითა და GitHub Pages-ით შეიქმნა. მას აქვს ინტერაქტიული ვიზუალიზაცია (Shiny, Plotly) და მომხმარებლებს შეუძლიათ
GitHub issues-ით ახალი წყაროები შესთავაზონ.

%==========================================================
\section{Data Records (მონაცემთა ჩანაწერები)} \pg{7--8}

უახლესი ნაკრები (სტატიის დროს ვერსია 1.2) ხელმისაწვდომია \href{https://opendengue.org/data.html}{opendengue.org}-ზე, ძველი და
ახალი ვერსიები კი GitHub-ზე (\href{https://github.com/OpenDengue/master-repo}{OpenDengue/master-repo}). ფაილები Figshare-ზეცაა
\texttt{csv} ფორმატში. მონაცემები \imp{FAIR სტანდარტებს} შეესაბამება.

სხვადასხვა რეზოლუციის მონაცემების გამო ავტორები სამ გლობალურ ამონაწერს გვთავაზობენ: საუკეთესო ეროვნული შეფასება, საუკეთესო
დროითი რეზოლუცია და საუკეთესო სივრცითი რეზოლუცია. \imp{ასე მომხმარებელს შეუძლია ამონაწერი თავისი კვლევითი კითხვის მიხედვით აირჩიოს.}

\textbf{სვეტები (codebook):} \texttt{adm\_0\_name}, \texttt{adm\_1\_name}, \texttt{adm\_2\_name}, \texttt{full\_name}, \texttt{ISO\_A0},
\texttt{FAO\_GAUL\_code}, \texttt{RNE\_iso\_code}, \texttt{IBGE\_code}, \texttt{calendar\_start\_date}, \texttt{calendar\_end\_date},
\texttt{Year}, \texttt{dengue\_total}, \texttt{case\_definition\_standardised}, \texttt{S\_res} (სივრცითი რეზოლუცია), \texttt{T\_res}
(დროითი რეზოლუცია), \texttt{UUID}.

\textbf{შეჯამება.} 56 მილიონზე მეტი შემთხვევა, 102 ქვეყანა, 1924--2023, \imp{843 სხვადასხვა წყარო}; ჩანაწერების 99.8\% კვირეული ან
თვიურია; სუბნაციონალური მონაცემები 40 ქვეყანაზე. მონაცემების 95\%-ზე მეტი სამინისტროებიდანაა. მომავალი პრიორიტეტები: წყნარი
ოკეანის კუნძულები (2011--2016), \imp{მეტი კვირეული მონაცემი აზიისთვის}, მეტი სუბნაციონალური მონაცემი სამხრეთ აზიისთვის.
ყველაზე მეტი შემთხვევა: ბრაზილია (22.0 მლნ), ვიეტნამი (4.5 მლნ), ფილიპინები (4.4 მლნ), ინდონეზია (2.8 მლნ).

\begin{etl}
\textbf{„843 სხვადასხვა წყარო“}: ეს რიცხვი კარგად აჩვენებს, რამდენად გაფანტულია ჯანმრთელობის მონაცემები. \textbf{„მეტი კვირეული
მონაცემი აზიისთვის“}: ანუ ტაილანდისთვის (აზია) კვირეული მონაცემები შეიძლება ნაკლები იყოს, ვიდრე ამერიკისთვის.
\end{etl}

%==========================================================
\section{Technical Validation (ტექნიკური ვალიდაცია) \textit{[შემოკლებით]}} \pg{8--9}

PAHO-ს მონაცემებში კვირეული ჯამები წლიურ ჯამებს შეადარეს. „Zero filling“-მა 52-დან 36 ქვეყნისთვის ხარვეზები შეავსო და
გამოტოვებული მონაცემების წილი 35.9\%-დან 28.9\%-მდე შეამცირა. სამმა ავტორმა დამოუკიდებლად შეამოწმა დროითი სერიების გრაფიკები
უჩვეულო ნახტომებზე და თარიღის შეცდომებზე. ყველა UUID წყაროს ცხრილს შეუდარდა. \imp{ვერსია 1.1-ში შეცდომამ ამერიკისთვის ზოგი
დუბლიკატი გამოიწვია, რაც 1.2-ში გასწორდა.} რეგულარულად ხდებოდა შემთხვევით შერჩეული სტრიქონების წყარომდე მიკვლევა.

\begin{etl}
\textbf{ვალიდაციის ეს მეთოდები შენს Case Study-საც გამოადგება}: (1) ჯამების შედარება (თვიური ჯამი = წლიური), (2) გრაფიკების
ვიზუალური შემოწმება, (3) შემთხვევითი სტრიქონის წყაროსთან შედარება. შეგიძლია თქვა, რომ შენ ვალიდაციის იგივე პრინციპებს იყენებ,
რასაც OpenDengue. ეს ასევე გვიჩვენებს, რომ \imp{ვერსიებს შორის მონაცემები იცვლება} (1.1 $\rightarrow$ 1.2), ამიტომ ვერსიის მითითება აუცილებელია.
\end{etl}

%==========================================================
\section{Usage Notes (გამოყენების შენიშვნები): სრული თარგმანი} \pg{9--12}

OpenDengue სხვადასხვა წყაროს მონაცემებს აერთიანებს და სტანდარტიზებს, რაც ახალ ანალიზებს გლობალურ და რეგიონულ დონეზე შესაძლებელს
ხდის. მაგალითებია: ყველაზე დაზარალებული რეგიონებისა და წლების გამოვლენა, \imp{გადაცემის მამოძრავებელი ფაქტორების, მაგალითად
კლიმატური ფაქტორებისა და ინტერვენციების, გაგება}, მომავალი ტენდენციებისა და აფეთქების რისკის პროგნოზირება.

\begin{citat}
\textit{„OpenDengue draws together and standardises data from multiple sources that enable new analyses at global and regional scales. Examples include identifying worst affected areas and years, understanding drivers of transmission such as climate factors and interventions and predicting future trends and outbreak risk.“} (Usage Notes, p.~9)
\end{citat}

\textbf{ამონაწერის არჩევა.} მომხმარებელმა უნდა გაითვალისწინოს, რომელი ამონაწერია (ეროვნული, სივრცითი თუ დროითი) ყველაზე
შესაფერისი მისი კვლევითი კითხვისთვის. \imp{დენგეს დინამიკის დროში ცვლილების შემსწავლელმა ანალიზებმა, მათ შორის დროითი სერიების
ანალიზმა, პროგნოზირებამ და ეროვნული პროგრამების შეფასებამ, დროითი ამონაწერი უნდა გამოიყენონ}, რომელიც მაქსიმალური დროითი
რეზოლუციის ჩანაწერებს ირჩევს. კონკრეტულ სუბნაციონალურ ადგილებზე ან გეოსივრცით ანალიზზე ორიენტირებულმა კვლევებმა სივრცითი
ამონაწერი უნდა გამოიყენონ. ეროვნული ამონაწერი თითოეული ქვეყნისთვის წლიური შემთხვევების ერთ საუკეთესო (უმაღლეს) შეფასებას იძლევა
და საუკეთესოდ ერგება ტვირთის შეფასებასა და ეროვნულ დონეზე გრძელვადიანი ტენდენციების ანალიზს. სამივე ამონაწერი წყაროებითა და
რიცხვებით დიდწილად ემთხვევა ერთმანეთს, მაგრამ ზოგ შემთხვევაში ზუსტ არჩევანს მნიშვნელობა აქვს.

\begin{citat}
\textit{„Applications that explore changes in dengue dynamics over time including time series analysis, forecasting and national programme evaluation should use the temporal extract which preferentially selects data records that maximise the temporal resolution of the data.“} (Usage Notes > Choice of extract, p.~10)
\end{citat}

\textbf{შედარებადობის შეზღუდვები.} როგორც ყველა ზედამხედველობის ბაზას, აქაც მნიშვნელოვანი შეზღუდვებია მონაცემთა შეგროვებისა და
ანგარიშგების სიზუსტესთან დაკავშირებით. \imp{დენგეს ზედამხედველობა, ანგარიშგება და მონაცემთა ხელმისაწვდომობა შეიძლება მნიშვნელოვნად
განსხვავდებოდეს ქვეყნებს შორის, ქვეყნის შიგნით, წყაროებს შორის და დროში}, რაც შედარებადობაზე მოქმედებს. ამიტომ სიფრთხილეს
გირჩევთ ქვეყნების შედარებისას ან ერთი ქვეყნის შიგნით დროში შედარებისას. შედარებადობის შესახებ ინფორმაციას სტანდარტიზებული შემთხვევის
განსაზღვრების ცვლადი იძლევა: „დადასტურებული“ შემთხვევების ჩანაწერები ნაკლებად მგრძნობიარე, მაგრამ უფრო სპეციფიკურია, ვიდრე
„სავარაუდო“ ან „საეჭვო“. თუმცა სტანდარტიზებული განსაზღვრებების შიგნითაც ეროვნული განსაზღვრებები მრავალფეროვანია, განსაკუთრებით
WHO-ს 1997 და 2009 წლის კრიტერიუმებს შორის გადასვლისას, რომლებიც ზოგჯერ 1997 ან 2009 წლის შემდეგ დიდი ხნით გვიან ამოქმედდა.
მომხმარებლებს ვურჩევთ, წყაროს მონაცემები და შემთხვევის განსაზღვრება შეამოწმონ. ყველა განსაზღვრებიდან გამოირიცხა შემოტანილი
შემთხვევები, სადაც ისინი ადგილობრივისგან ცალკე იყო მოცემული, მაგრამ მათი გარჩევის მეთოდები განსხვავებულია. გადაცემის
ჩრდილოეთ და სამხრეთ საზღვრებთან ზოგი მონაცემი შეიძლება შერეულ შემთხვევებს შეიცავდეს. ასევე მნიშვნელოვანია ითქვას, რომ \imp{კარგი
ზედამხედველობის სისტემის მქონე ქვეყნებშიც კი აღრიცხული შემთხვევები ყველა ინფექციის მცირე და ცვალებად წილს შეადგენს}: უსიმპტომო
ინფექციის, მკურნალობაზე მიმართვის განსხვავებების, კერძო სექტორში მკურნალობისა და პირველადი ჯანდაცვის დიაგნოსტიკის სირთულეების გამო.

კონკრეტული ქვეყნების ან რეგიონების დეტალური ანალიზისთვის მომხმარებლებს ადგილობრივ ექსპერტებთან ან ჯანდაცვის უწყებებთან
დაკავშირებას ვურჩევთ. ეს დაეხმარება იმის გაგებას, როგორ იქმნება მონაცემები, და იმის შემოწმებას, ხომ არ არის დაკვირვებული
ნიმუშები ან დასკვნები შემთხვევის განსაზღვრების ან ანგარიშგების პრაქტიკის ცვლილების შედეგი.

\begin{etl}
\textbf{შენი სისტემისთვის:} ეს შეზღუდვები შენს მონაცემებზეც გადადის. LLM-ის ახსნაში ან PDF ანგარიშში შეიძლება გაფრთხილება დაემატოს:
„შემთხვევის განსაზღვრება და ანგარიშგება შეიძლება დროში შეიცვალოს, ამიტომ ქვეყნებს შორის შედარება სიფრთხილით გააკეთეთ“.
თუ სისტემა \texttt{case\_definition\_standardised}-ს შედეგში გადაიტანს, ეს მეტ გამჭვირვალობას მისცემს (რეკომენდაცია, P2).
\end{etl}

\textbf{ვერსიების კონტროლი.} OpenDengue მუდმივად ვითარდება და ახალი ვერსიები სულ მცირე ექვს თვეში ერთხელ გამოდის, Figshare-ზე
სხვადასხვა DOI-ით. \imp{მომხმარებლებს ვურჩევთ, მიუთითონ, OpenDengue-ს რომელი ვერსია გამოიყენეს}, და პროექტის მნიშვნელოვან
ეტაპებზე განახლებები შეამოწმონ. მომავალ ვერსიებში დაგეგმილია მონაცემების დაყოფა სიმძიმის, დადასტურების მეთოდის, ასაკისა და
სეროტიპის მიხედვით.

\begin{citat}
\textit{„It is recommended for users to specify which version of OpenDengue they use in their analyses and routinely check for updates at relevant points in their project lifecycle.“} (Usage Notes > Dataset version control, p.~12)
\end{citat}

\textbf{ციტირება და ლიცენზია.} \imp{OpenDengue-ს მონაცემები Creative Commons CC BY-SA ლიცენზიით ვრცელდება.} ეს ყველა მომხმარებელს
(კომერციულსა და არაკომერციულს) აძლევს უფლებას, მონაცემები სათანადო მითითებით ხელახლა გამოიყენოს და შეცვალოს.
\imp{CC BY-SA ლიცენზიით OpenDengue-ს ნაკრების ყველა ადაპტაცია იმავე პირობებით უნდა გავრცელდეს.} სასურველი ციტირებაა ეს სტატია
და გამოყენებული ვერსიის Figshare-ის ბმული:

\begin{quote}\small
Clarke, Joe; Lim, Ahyoung; Gupte, Pratik R.; Pigott, David M.; van Panhuis, Wilbert G; Brady, Oliver (2023). OpenDengue: data from the
OpenDengue database. Version [1.2]. figshare. Dataset. \url{https://doi.org/10.6084/m9.figshare.24259573}
\end{quote}
შესაძლებლობის შემთხვევაში მომხმარებლებს ორიგინალი წყაროების ციტირებასაც ვურჩევთ (UUID-ით).

\begin{alarm}
\textbf{ლიცენზია CC BY-SA (ShareAlike) შენს სისტემას პირდაპირ ეხება!}
\begin{itemize}[leftmargin=1.5em]
  \item შენი სისტემის გამოსატანი ცხრილი, რომელიც OpenDengue-ს მონაცემებს შეიცავს, \textbf{ადაპტაციაა}. ამიტომ ის
        \imp{იმავე CC BY-SA ლიცენზიით} უნდა გავრცელდეს და OpenDengue მითითებული უნდა იყოს.
  \item ახლა სისტემა სტატიას ციტირებს (\texttt{config.py}), მაგრამ \textbf{არც მონაცემთა ნაკრების ვერსიას} (Figshare DOI) და \textbf{არც
        ლიცენზიას} არ აჩვენებს.
  \item \textbf{რეკომენდაცია:} PDF ანგარიშსა და CSV-ის მეტამონაცემებში ჩაწერე: OpenDengue ვერსია + Figshare DOI + „CC BY-SA“ (ვერსია: შეამოწმე opendengue.org-ზე).
        ეს თეზისშიც უნდა ახსენო (სამართლებრივი და ეთიკური ასპექტი: \textit{Lizenzierung}).
\end{itemize}
\end{alarm}

\textbf{მონაცემების შეთავაზება და უკუკავშირი.} თუ მომხმარებლებმა იციან მონაცემების შესახებ, რომლებიც ბაზაში აკლია, მათი
გაზიარება მისასალმებელია (\href{https://opendengue.org/contribute}{opendengue.org/contribute}). შეცდომების შეტყობინებაც მისასალმებელია.

\textbf{კოდის ხელმისაწვდომობა.} მონაცემების დამუშავებისა და სტანდარტიზაციის მთელი კოდი OpenDengue-ს GitHub საცავშია.

%==========================================================
\section{შეჯამება შენთვის}

\subsection*{ციტატები თეზისისთვის (მოტივაცია)}
\begin{longtable}{@{}p{0.6\linewidth}p{0.12\linewidth}p{0.22\linewidth}@{}}
\toprule
\textbf{ციტატა (შემოკლებით)} & \textbf{გვერდი} & \textbf{რას ამტკიცებს}\\
\midrule
„\ldots difficulties in collecting and aggregating consistent and comparable dengue incidence\ldots data“ & 1 & მონაცემების შეგროვება რთულია\\
„they can often be difficult to find and no single database aggregates data\ldots“ & 1 & მონაცემები გაფანტულია\\
„\ldots better understand the drivers of transmission\ldots model\ldots under different climate scenarios“ & 1 & კლიმატთან გაერთიანების საჭიროება\\
„\ldots requires substantial manual processing“ (ProMED) & 2 & ხელით დამუშავება\\
„Locating, extracting, processing and standardising data all takes time\ldots“ & 2 & \textbf{მთავარი ციტატა}: ETL შრომატევადია\\
„\ldots traced back to their original reporting source\ldots“ & 2 & provenance (FR-18)\\
„\ldots understanding drivers of transmission such as climate factors\ldots“ & 9 & გამოყენების მიზანი = შენი სისტემა\\
„\ldots should use the temporal extract\ldots“ & 10 & მოთხოვნა: სწორი ამონაწერი\\
„\ldots specify which version of OpenDengue they use\ldots“ & 12 & მოთხოვნა: ვერსიის მითითება\\
\bottomrule
\end{longtable}

\subsection*{რა შეამოწმო ან დაამატო შენს სისტემაში}
\begin{longtable}{@{}p{0.05\linewidth}p{0.55\linewidth}p{0.08\linewidth}p{0.26\linewidth}@{}}
\toprule
\textbf{\#} & \textbf{საკითხი} & \textbf{პრ.} & \textbf{წყარო სტატიაში}\\
\midrule
1 & რომელი \textbf{ამონაწერიდან} (temporal / national / spatial) და რომელი \textbf{ვერსიიდან} შეიქმნა \texttt{dengue\_monthly\_by\_iso3.csv}? დოკუმენტაციაში ჩაწერე & P0 & p.~10, 12\\
2 & \textbf{ლიცენზია CC BY-SA} + Figshare DOI + ვერსია შედეგში და PDF-ში & P0 & p.~12\\
3 & როგორ გადავიდა \textbf{კვირეული $\rightarrow$ თვიური} (ორ თვეზე გადასული ეპი-კვირები)? & P1 & p.~5\\
4 & ტაილანდის მონაცემები ფაილში \textbf{2009-01-დან} იწყება, Suttirat კი 2003-დან. შეამოწმე, აქვს თუ არა OpenDengue-ს ადრეული მონაცემები & P1 & ---\\
5 & დამატებითი ან იმპუტირებული მნიშვნელობების \textbf{მონიშვნა} შედეგში & P2 & p.~5--7\\
6 & \texttt{case\_definition\_standardised} შედეგში + შედარებადობის გაფრთხილება & P2 & p.~10--12\\
\bottomrule
\end{longtable}

\subsection*{Related Work-ისთვის}
სტატიის \textbf{Table 1} (გვ.~2) ჯანმრთელობის მონაცემთა ბაზების მზა შედარებაა: Project Tycho, PAHO PLISA, ECDC, GIDEON, ProMED.
ის შეგიძლია ციტირო, როცა აჩვენებ, რომ \textbf{ჯანმრთელობის მხარეს} უკვე არსებობს კარგი წყაროები, მაგრამ არცერთი მათგანი
\textbf{კლიმატთან გაერთიანებას} არ აკეთებს.

\end{document}
